% TOP_PROBLEM: {"id": 3059, "title": "A nowhere-zero point in a linear mapping", "category_id": 2, "category": {"id": 2, "name": "combinatorics", "display_name": "Combinatorics", "description": "Counting problems, graph theory, discrete structures.", "slug": "combinatorics", "order_index": 2, "created_at": "2026-07-31T15:26:25.671Z"}, "status": "open", "problem_number": "OPG-148", "set_id": 12}
% FIRST_POSTED: 2026-10-01
% ENTRY_KIND: statement_only
% UnsolvedMath ID 3059; source catalogue code OPG-148
% !TeX program = lualatex
% Definition and sources only; this file is not an LLM solution attempt.
\documentclass[11pt,a4paper]{article}
\usepackage[margin=25mm]{geometry}
\usepackage{fontspec}
\setmainfont{lmroman10-regular.otf}[BoldFont=lmroman10-bold.otf,
 ItalicFont=lmroman10-italic.otf,BoldItalicFont=lmroman10-bolditalic.otf]
\usepackage{amsmath,amssymb,mathtools}
\usepackage{unicode-math}
\setmathfont{latinmodern-math.otf}
\AtBeginDocument{\let\setminus\smallsetminus}
\usepackage{xurl,hyperref}
\hypersetup{colorlinks=true,urlcolor=blue}
\setcounter{secnumdepth}{0}
\setlength{\parindent}{0pt}
\setlength{\parskip}{5pt}
\setlength{\emergencystretch}{3em}
\begin{document}
\section*{A nowhere-zero point in a linear mapping}\label{3059}
{\small UnsolvedMath ID 3059 · OPG-148 · Combinatorics · OpenGarden}
\subsection{Definitions and mathematical statement}
Let $\mathbb F_q$ be a finite field of cardinality $q\ge4$, let $n\ge1$,
and let $A\in\mathbb F_q^{n\times n}$ be invertible, meaning
$\det(A)\ne0$ in the field. Put
$\mathbb F_q^\times=\mathbb F_q\setminus\{0\}$ and
$T=(\mathbb F_q^\times)^n$.
A vector is nowhere zero when each of its coordinates is nonzero; this
is stronger than requiring the vector itself to differ from the zero vector.

The conjecture states
\[
 \forall q\ge4\text{ a prime power}\quad\forall n\ge1\quad
 \forall A\in\operatorname{GL}_n(\mathbb F_q),\qquad
 A(T)\cap T\ne\varnothing.
\]
Equivalently, there are $x_1,\ldots,x_n\in\mathbb F_q^\times$ such that
$\sum_{j=1}^n a_{ij}x_j\ne0$ for every row $i$. Setting $y=Ax$ then
supplies the two nowhere-zero vectors appearing in the original statement.
Both vectors are over the same field; their coordinates may repeat.

Another description asks whether the $n$ coordinate hyperplanes
$x_j=0$, together with the $n$ hyperplanes
$\sum_j a_{ij}x_j=0$, can ever cover all of $\mathbb F_q^n$ when the rows
of $A$ are independent. The desired answer is that they cannot.
The matrix and dimension are arbitrary, and any existence argument must
work without additional assumptions on the pattern of nonzero entries.
The lower bound $q\ge4$ is part of the assertion and is retained.
\subsection{Short English statement}
Must every invertible matrix over a finite field with at least four elements send some coordinatewise nonzero vector to another such vector?
\subsection{Sources}
\begin{itemize}
\item[S1] ULAM.AI and the UnsolvedMath Contributors, supplied problems.json, record 3059 (OPG-148), OpenGarden collection. Catalogue identity and source transcription; CC BY 4.0.
\url{https://huggingface.co/datasets/ulamai/UnsolvedMath}.
\item[S2] Open Problem Garden, A nowhere-zero point in a linear mapping. Original problem and discussion; the mathematical formulation below is expanded from the supplied source record.
\url{http://www.openproblemgarden.org/op/a_nowhere_zero_point_in_a_linear_mapping}.
\end{itemize}
\end{document}
